\chapter{Discussion}
\label{Discussion}

This chapter interprets the experimental results of Chapter~\ref{AER} in the context of the contributions outlined in the Introduction and the prior work reviewed in Chapter~\ref{LR}. We organise the discussion around three themes: the predictive performance of the hybrid model and its architectural implications, the interpretability of the learned closure as revealed by the spectral gain analysis, and the limitations that qualify the conclusions.

\section{Hybrid architecture: predictive performance and data efficiency}

The Introduction motivated the hybrid formulation as a middle ground between purely numerical and purely data-driven approaches, claiming benefits in accuracy, stability, and data efficiency. The ablation results in Table~\ref{tab:ablation_ml_dycore_only} support this claim directly: the hybrid model (spectral dycore + GNN closure) outperforms both the dycore-only and the ML-only variants at the same lead time and resolution. Notably, the hybrid variant trained on $\sim\!1900$ data points surpasses the ML-only model trained on $12$K points, demonstrating the data efficiency afforded by encoding physical structure in the solver. This is consistent with the broader argument in the hybrid-modelling literature that mechanistic inductive biases reduce the sample complexity of the learning problem \citep{Kashinath2021, Meng2025, Kochkov2021}.

Equally important is the ablation comparing online (trajectory-fitting) and offline (single-step) training (Figure~\ref{fig:rmse_horizon}). Online training improves performance not only at long lead times, where distributional shift is expected, but already at the shortest 6\,h horizon. This confirms the theoretical motivation in Section~\ref{sec:formal_closure}: because online training optimises the closure on the distribution of states actually produced by the coupled solver, it reduces the mismatch between training and deployment that plagues offline closures \citep{kurz2021investigating, brenowitz2019spatially}. The widening gap between online and offline RMSE with increasing lead time further underlines the importance of differentiable solvers, which enable backpropagation through the full simulation pipeline, a design choice that was central to our model architecture.

\section{Resolution--performance trade-offs}

A recurring finding in Chapter~\ref{AER} is the non-monotonic relationship between grid resolution and forecast skill. At short lead times, finer grids generally improve MSE, consistent with the expectation that higher resolution recovers additional dynamical variability. However, at longer lead times this ordering inverts: TL31 or TL63 often yield lower MSE than TL95 or TL127. We attribute this to the interaction of three mechanisms.

First, chaotic error growth at fine scales causes high-wavenumber structures to decorrelate rapidly, so that small initial errors amplify with rollout horizon. Second, the resolution-aware GNN scaling (Section~\ref{sec:gnn_scaling}), which reduces per-node hidden width as the number of nodes increases in order to keep compute constant, progressively limits the expressive capacity of the closure at finer grids. Third, MSE is subject to the well-known double-penalty effect: a sharp structure predicted with correct amplitude but slight spatial displacement incurs large squared error, whereas a smoother, coarser field avoids this penalty \citep{mass2002does}. Together, these effects explain why TL63 frequently acts as a practical sweet spot for dynamical variables, retaining sufficient spatial detail while avoiding the instability and capacity stress observed at TL127.

The variable-dependent character of this trade-off is also notable. Divergence strongly favours coarse resolution, suggesting that it is particularly sensitive to high-wavenumber noise. In contrast, specific humidity benefits from finer grids up to TL95, consistent with the localised, small-scale nature of moisture dynamics, before degrading again at TL127 due to capacity dilution. These contrasts argue against selecting a single grid resolution on the basis of one aggregate score and in favour of jointly tuning resolution, closure capacity, and training objective for the downstream target.

\section{Interpreting the learned closure through the gain analysis}
\label{sec:discuss_gain}

The spectral gain experiment, formalised in Section~\ref{sec:gain-experiment}, constitutes the primary interpretability contribution of this thesis. By linearising the learned closure around representative coupled states and projecting the resulting Jacobian--tendency product onto the spherical-harmonic basis, we obtain a quantitative, scale-resolved picture of how the closure modifies the dynamical core's tendencies. In this section we discuss the principal findings and relate them to the interpretability literature reviewed in Section~\ref{LR}.

\subsection{Predominantly amplifying character of the closure}

Across resolutions and lead times, the gain $\bar{G}_\ell$ is above unity for a majority of wavenumbers, indicating that the closure predominantly reinforces, rather than damps, the explicit tendencies of the dynamical core. In other words, the learned component provides a net positive feedback that partially counteracts the aggregate damping introduced by spectral truncation, filtering, and implicit numerical diffusion. This finding is consistent with the energy-backscatter mechanism identified by \citet{guan2022stable}, who showed that a neural closure in two-dimensional turbulence simultaneously applies forward diffusion for stability and backscatters energy to restore missing subgrid variance. In our setting the backscatter analogue manifests as tendency amplification, while the stabilisation manifests as the damping fraction that grows with lead time.

The result also aligns with the observations of \citet{subel2023explaining}, whose SPARK analysis of convolutional closures in two-dimensional turbulence revealed that network layers function as explicit spectral filters that bridge the energy gap between resolved and subgrid scales. Our Jacobian-based gain plays a conceptually similar role---it quantifies the spectral redistribution of tendency energy by the closure---but extends the analysis to a three-dimensional, multi-variable atmospheric model operating on the sphere, where the spherical-harmonic basis replaces the Cartesian Fourier basis.

\subsection{Scale selectivity: damping at large scales, amplification at small scales}

A key finding is that the closure's behaviour is not uniform across scales. At low wavenumbers (large spatial scales), the fraction of damping modes $f^{-}_\ell$ is highest, while at high wavenumbers (small spatial scales) $f^{+}_\ell$ dominates. This scale-selective character suggests a physically interpretable division of labour: the closure stabilises the slowly evolving, large-scale components of the flow, while restoring small-scale variability that would otherwise be lost to truncation and filtering.

This interpretation connects to two strands in the literature. First, it mirrors the dual role reported by \citet{Portwood2021}, who computed exact Jacobian-based sensitivities of a neural closure with respect to localised flow gradients in isotropic turbulence and found that the network dynamically balances dissipation and structural reconstruction, mirroring classical mixed-model theories. Our spectral gain extends this Jacobian-based perspective to a global, multi-scale setting and confirms that a qualitatively similar balance emerges even in a vastly more complex system. Second, the finding that the closure acts as a scale-dependent energy redistributor is consistent with the eddy-closure diagnostics of \citet{ross2023benchmarking}, who found that neural closures maintain long-term spectral energy fidelity more effectively than symbolic closures, precisely because they can modulate backscatter and forward transfer across a wide range of scales. Our gain analysis provides a mechanistic explanation: the closure achieves this fidelity by selectively amplifying tendency energy at under-resolved scales and damping it at well-resolved ones.

\subsection{Signature of the spectral filter in the gain spectrum}

A particularly striking feature is the pronounced increase in $f^{+}_\ell$ near the eighth decile of representable wavenumbers, accompanied by a corresponding dip in $f^{-}_\ell$. This coincides with the onset of the exponential spectral filter (Section~\ref{subsec:filtering}), which explicitly attenuates high-wavenumber modes. The closure has learned to partially compensate for this artificial damping by amplifying tendency energy at the affected scales.

This observation has two important implications. First, it demonstrates that the gain diagnostic can distinguish between physical and numerical contributions to the closure's behaviour. The filter-compensation signature is a numerical artefact of the model design, not a reflection of missing subgrid physics. The ability to identify such artefacts is a practical benefit of the proposed diagnostic that goes beyond what aggregate forecast scores can reveal. Second, the renewed increase in damping modes very close to the truncation limit, which we attribute to aliasing and transform errors near the highest retained modes, illustrates that even a well-trained closure is affected by the discretisation environment in which it operates. Both of these points underscore the value of scale-resolved diagnostics over black-box evaluation.

\subsection{Lead-time dependence and the transition from amplification to damping}

The gain analysis reveals a systematic shift with forecast horizon: as lead time increases, the fraction of amplifying modes decreases while the fraction of damping modes increases. This finding connects to the well-known over-smoothing tendency in learned chaotic dynamical systems \citep{pathak2022fourcastnet}. As the rollout horizon grows, small errors amplify due to chaotic dynamics, and the MSE training objective increasingly penalises sharp, energetically active structures. The closure responds by becoming progressively more dissipative. In the gain framework, this manifests as a shift of $\bar{G}_\ell$ toward and below unity.

This transition is not linear: the $f^{+}$ and $f^{-}$ curves exhibit a knee-like, piecewise structure whose inflection point depends on grid resolution. Since the rollout-length checkpoints during training were identical across resolutions, this resolution dependence is unlikely to be a direct artefact of the training schedule. Instead, it may reflect the resolution-dependent decorrelation time of high-wavenumber modes: at finer grids, chaotic decorrelation sets in earlier, causing the closure to shift to a damping-dominated regime at shorter lead times.

These observations extend the findings of \citet{guan2022stable}, who documented a similar balance between forward dissipation and energy backscatter in their two-dimensional turbulence closure, but only at a single, fixed integration time. Our analysis shows that this balance is dynamic: the relative strength of amplification and damping shifts systematically with the forecast horizon, and does so in a resolution-dependent manner.

\subsection{Lower resolution and stronger dissipation}

Across all diagnostics, lower-resolution grids exhibit a larger fraction of damping modes, especially at long lead times and low normalised wavenumbers. This is consistent with the stronger effective dissipation and reduced scale separation at coarse resolution, where truncated nonlinear interactions and numerical filtering damp variability more aggressively. Consequently, the closure at coarse resolution has less under-resolved variability to restore and instead acts more diffusively to maintain rollout stability. At higher resolution, the closure retains a larger share of amplifying modes, reflecting a richer spectrum of under-resolved processes that the network can partially compensate.

\section{Situating the gain diagnostic in the interpretability landscape}

A central argument of this thesis is that the GNN-based closure architecture, combined with the fully differentiable simulation pipeline, enables a form of empirical interpretability that was not previously available for realistic hybrid atmospheric models. It is instructive to compare the proposed gain diagnostic with the interpretability methods reviewed in Chapter~\ref{LR}.

\citet{subel2023explaining} analysed convolutional closures by transforming network \emph{weights} into the frequency domain (SPARK), which reveals the spectral character of each learned filter. This approach is architecturally specific: it relies on the translational structure of convolutions to interpret weight matrices as spatial filters. Our gain analysis, by contrast, operates on the input--output behaviour of the closure (via the Jacobian) rather than on its internal weights, making it agnostic to the specific architecture of the learned component. This is important because it means the diagnostic is, in principle, transferable to any differentiable closure---a point enabled by our choice to implement the entire pipeline in JAX.

\citet{Portwood2021} also used Jacobian-based analysis, computing sensitivities of a neural closure with respect to local flow gradients. Their approach shares the same mathematical foundation as ours (linearisation of the learned operator), but was applied in a single-scale isotropic turbulence setting without projecting the Jacobian onto a spectral basis. By projecting the Jacobian--tendency product onto the spherical-harmonic basis, we extend their approach to a multi-scale diagnostic that resolves the closure's effect at each spatial wavenumber.

The symbolic-regression approach of \citet{zanna2020data} offers a complementary form of interpretability: closed-form algebraic expressions that can be inspected directly for physical content. However, as \citet{ross2023benchmarking} demonstrated, symbolic closures sacrifice predictive fidelity relative to high-capacity neural networks. Our approach occupies a different point in this trade-off: we retain a black-box neural network for predictive power but interrogate its behaviour through external diagnostics. The gain analysis does not reveal the internal mechanism of the network in the way a symbolic expression does, but it does reveal the \emph{functional effect} of the closure at each scale and lead time, which is often the quantity of practical interest.

Finally, the choice of a GNN over NeuralGCM's column-wise MLP \citep{neuralgcm} is critical for the interpretability programme. Because NeuralGCM's MLP operates independently on each vertical column, it is fundamentally unable to express or reveal horizontal scale selectivity. By implementing the closure as a spatially connected GNN, our architecture enables the Jacobian to carry horizontal spatial structure, which is what makes the wavenumber-resolved gain analysis possible. The graph Laplacian eigenmodes of the GNN mesh provide a natural analogue to spherical harmonics, making the proposed cross-scale analysis adaptable to other settings where a GNN serves as the closure operator.

\section{Limitations}

Several limitations qualify the conclusions drawn above.

\begin{itemize}
    \item \textbf{Reliance on MSE as training and evaluation metric.} Both the training loss and the primary evaluation metric are based on mean squared error, which penalises phase errors disproportionately and may under-reward physically plausible but spatially displaced structures. This affects both the forecast skill comparisons (through the double-penalty effect) and, indirectly, the gain analysis, since the closure is optimised under this objective. Alternative loss functions, such as spectral or perceptual losses, might lead to qualitatively different closure behaviours.
    
    \item \textbf{Linearisation assumption in the gain diagnostic.} The gain analysis is based on a first-order Taylor expansion of the learned closure. While this linearisation is exact for infinitesimal perturbations, the actual tendency magnitudes are finite, so nonlinear effects (higher-order terms in the Taylor expansion) are not captured. The diagnostic therefore characterises the closure's local, linear interaction with the resolved dynamics, which may not fully reflect its behaviour in strongly nonlinear regimes.
    
    \item \textbf{Single prognostic variable for the gain.} The gain results are computed using the temperature-variation field as the representative prognostic variable. While temperature is a natural choice given its central role in the primitive equations, other variables (e.g.\ vorticity, divergence, moisture) may exhibit different gain characteristics. Extending the analysis to all prognostic variables would provide a more complete picture but was beyond the scope of this work.
    
    \item \textbf{Finite test period.} Gain statistics and rollout evaluations are computed on a single-year test window (2020), which may not span rare or extreme regimes. The ensemble-averaged gain $\bar{G}_\ell$ suppresses flow-specific fluctuations but may not generalise to markedly different climate conditions.
    
    \item \textbf{Capacity--compute confound at high resolution.} The resolution-aware GNN scaling keeps compute comparable across grids but confounds ``more resolution'' with ``less per-node model capacity''. The performance degradation at TL127 may therefore reflect under-parameterisation of the closure rather than a fundamental resolution effect. Disentangling these two factors would require experiments with matched per-node capacity, at correspondingly higher computational cost.
    
    \item \textbf{Deterministic evaluation only.} All results focus on deterministic trajectories and ensemble-mean diagnostics. The gain analysis does not address how the closure affects predictive uncertainty or ensemble spread, which are important considerations for operational forecasting.
    
    \item \textbf{Stability ceiling.} NaN and exploding outcomes at selected long-horizon settings indicate that robustness is not yet uniform across all resolutions and lead times, limiting the practical forecast range of the current model.
\end{itemize}

These limitations do not invalidate the core findings, but they indicate that some conclusions are conditional on the metric, training design, and operational horizon chosen. They also point toward natural extensions of this work, which are discussed in Chapter~\ref{CAFW}.
